%==============================================================================
%  Trabajo de Fin de Grado (TFG)
%  Análisis Topológico de Datos aplicado a la clasificación de mamografías
%  Plantilla base -- generada como punto de partida
%==============================================================================
\documentclass[12pt,a4paper,twoside,openright]{book}

%------------------------------- Codificación e idioma ------------------------
\usepackage[utf8]{inputenc}
\usepackage[T1]{fontenc}
\usepackage[spanish,es-tabla]{babel}
\usepackage{lmodern}

%------------------------------- Márgenes y diseño ----------------------------
\usepackage[a4paper,top=2.5cm,bottom=2.5cm,left=3cm,right=2.5cm]{geometry}
\usepackage{setspace}
\onehalfspacing
\usepackage{emptypage}   % páginas en blanco realmente vacías

%------------------------------- Gráficos y tablas ----------------------------
\usepackage{graphicx}
\usepackage{float}
\usepackage{booktabs}
\usepackage{multirow}
\usepackage{array}
\usepackage{caption}
\usepackage{subcaption}
\usepackage{xcolor}

%------------------------------- Matemáticas ----------------------------------
\usepackage{amsmath,amssymb,amsthm}
\usepackage{mathtools}

%------------------------------- Código ---------------------------------------
\usepackage{listings}
\lstset{
  basicstyle=\ttfamily\footnotesize,
  keywordstyle=\color{blue!70!black}\bfseries,
  commentstyle=\color{green!40!black},
  stringstyle=\color{red!60!black},
  numbers=left, numberstyle=\tiny\color{gray},
  frame=single, breaklines=true, showstringspaces=false,
  captionpos=b, language=Python
}

%------------------------------- Referencias ----------------------------------
\usepackage[hidelinks]{hyperref}
\usepackage{cleveref}
\usepackage[backend=biber,style=ieee,sorting=nyt]{biblatex}
\addbibresource{bibliografia.bib}

%------------------------------- Teoremas -------------------------------------
\theoremstyle{definition}
\newtheorem{definicion}{Definición}[chapter]
\newtheorem{teorema}{Teorema}[chapter]
\newtheorem{ejemplo}{Ejemplo}[chapter]

%------------------------------- Encabezados ----------------------------------
\usepackage{fancyhdr}
\pagestyle{fancy}
\fancyhf{}
\fancyhead[LE,RO]{\thepage}
\fancyhead[RE]{\leftmark}
\fancyhead[LO]{\rightmark}
\renewcommand{\headrulewidth}{0.4pt}

%==============================================================================
\begin{document}

%------------------------------- Portada --------------------------------------
%------------------------------- Portada --------------------------------------
\begin{titlepage}
\centering

% --- Logo de la universidad (sustituir por el fichero real) ---
% \includegraphics[width=0.35\textwidth]{figuras/logo_universidad.png}\\[1.5cm]

{\scshape\LARGE Universidad de Alicante \par}
{\scshape\large Escuela Politécnica Superior \par}
\vspace{0.5cm}
{\large Grado en Ingeniería Informática \par}   % <-- ajustar titulación
\vspace{2cm}

\rule{\linewidth}{0.5mm}\\[0.6cm]
{\huge\bfseries Análisis Topológico de Datos aplicado a la clasificación y visualización de mamografías digitales \par}
\vspace{0.4cm}
\rule{\linewidth}{0.5mm}\\[1.5cm]

{\large Trabajo de Fin de Grado \par}
\vspace{2.5cm}

\begin{minipage}{0.45\textwidth}
\begin{flushleft}\large
\textit{Autor:}\\
Nombre Apellido Apellido
\end{flushleft}
\end{minipage}
\hfill
\begin{minipage}{0.45\textwidth}
\begin{flushright}\large
\textit{Tutor/es:}\\
Nombre del tutor/a
\end{flushright}
\end{minipage}

\vfill
{\large \today \par}

\end{titlepage}


\frontmatter
\pagenumbering{roman}

%------------------------------- Resumen --------------------------------------
\chapter*{Resumen}
\addcontentsline{toc}{chapter}{Resumen}
% ----------------------------------------------------------------------------
% Resumen / Abstract
% ----------------------------------------------------------------------------

\noindent
El cáncer de mama es el tumor de mayor incidencia entre las mujeres a nivel
mundial, y la mamografía sigue siendo la prueba de cribado de referencia para su
detección precoz. En los últimos años, las técnicas de \textit{deep learning}
han demostrado un gran potencial en el análisis automático de imagen médica,
pero suelen comportarse como cajas negras y son sensibles al ruido y a la
variabilidad de adquisición. El \textbf{Análisis Topológico de Datos} (TDA, por
sus siglas en inglés) ofrece una perspectiva complementaria: cuantifica la
\emph{forma} de los datos mediante invariantes topológicos robustos frente al
ruido, como la homología persistente.

\medskip
\noindent
Este Trabajo de Fin de Grado estudia la aplicación del TDA a la clasificación y
visualización de mamografías digitales del conjunto de datos público
\mbox{CBIS-DDSM}. Se extraen descriptores topológicos (diagramas de
persistencia, códigos de barras e imágenes de persistencia) a partir de las
imágenes y de las regiones de interés, se evalúa su capacidad discriminativa
entre lesiones benignas y malignas, y se comparan y combinan con modelos de
aprendizaje profundo. El objetivo último es obtener un sistema de apoyo al
diagnóstico más \emph{interpretable} sin sacrificar rendimiento.

\vspace{1cm}
\noindent\textbf{Palabras clave:} análisis topológico de datos, homología
persistente, mamografía, cáncer de mama, aprendizaje profundo, imagen médica,
CBIS-DDSM.

\bigskip\bigskip
\begin{center}\rule{0.5\textwidth}{0.4pt}\end{center}
\bigskip

\selectlanguage{english}
\noindent\textbf{Abstract.}
Breast cancer is the most common malignancy among women worldwide, and
mammography remains the reference screening test for its early detection. Deep
learning has shown remarkable potential in medical image analysis, yet such
models often behave as black boxes and are sensitive to noise and acquisition
variability. \textbf{Topological Data Analysis} (TDA) provides a complementary
viewpoint by quantifying the \emph{shape} of the data through noise-robust
topological invariants such as persistent homology. This bachelor's thesis
investigates the application of TDA to the classification and visualisation of
digital mammograms from the public CBIS-DDSM dataset. Topological descriptors
(persistence diagrams, barcodes and persistence images) are extracted and
evaluated for their ability to discriminate benign from malignant lesions, and
compared and combined with deep learning models, aiming at a more
\emph{interpretable} decision-support system.

\medskip
\noindent\textbf{Keywords:} topological data analysis, persistent homology,
mammography, breast cancer, deep learning, medical imaging, CBIS-DDSM.
\selectlanguage{spanish}


%------------------------------- Índices --------------------------------------
\tableofcontents
\listoffigures
\listoftables

%------------------------------- Cuerpo ---------------------------------------
\mainmatter
\pagenumbering{arabic}

\chapter{Introducción}
\label{cap:introduccion}

\section{Motivación}
El cáncer de mama es la neoplasia más frecuente en mujeres y una de las
principales causas de mortalidad por cáncer a nivel mundial. La detección precoz
mediante programas de cribado con mamografía reduce de forma significativa la
mortalidad, ya que permite identificar lesiones en estadios tempranos y más
tratables. Sin embargo, la interpretación de mamografías es una tarea compleja,
sujeta a variabilidad inter-observador y a una carga de trabajo elevada para los
radiólogos.

En este contexto, los sistemas de diagnóstico asistido por computador
(\textit{Computer-Aided Diagnosis}, CAD) y, más recientemente, las técnicas de
aprendizaje profundo (\textit{deep learning}) se han convertido en herramientas
de gran interés. Trabajos como el de Yala \textit{et al.}~\cite{yala2019} han
demostrado que los modelos basados en mamografía pueden mejorar la predicción del
riesgo de cáncer de mama. No obstante, estos modelos presentan limitaciones
importantes: son difíciles de interpretar, requieren grandes cantidades de datos
etiquetados y pueden ser sensibles al ruido y a las condiciones de adquisición.

El \textbf{Análisis Topológico de Datos} (TDA) surge como un enfoque
complementario que caracteriza la \emph{forma} intrínseca de los datos mediante
herramientas de la topología algebraica. Su técnica principal, la
\emph{homología persistente}, es robusta frente al ruido y proporciona
descriptores interpretables geométricamente, lo que la hace especialmente
atractiva para el análisis de imagen médica.

\section{Contexto del problema}
Una mamografía digital es una imagen de rayos X de la mama que revela
estructuras como el tejido glandular, los conductos, las calcificaciones y
posibles masas. La densidad mamaria (clasificada habitualmente según el sistema
BI-RADS) y la morfología de las lesiones son factores clave en el diagnóstico.
El conjunto de datos público \mbox{CBIS-DDSM}~\cite{lee2017cbisddsm} proporciona
mamografías digitalizadas con anotaciones de lesiones (masas y
microcalcificaciones), regiones de interés (ROI) segmentadas y etiquetas de
patología (benigno/maligno), lo que lo convierte en un banco de pruebas
adecuado para este trabajo.

\section{Objetivos y alcance (resumen)}
El objetivo general de este TFG es \textbf{estudiar y evaluar la aplicación del
Análisis Topológico de Datos a la clasificación y visualización de mamografías},
comparándolo y combinándolo con técnicas de aprendizaje profundo. Los objetivos
específicos se detallan en el Capítulo~\ref{cap:objetivos}.

\section{Estructura de la memoria}
El resto de la memoria se organiza como sigue:
\begin{itemize}
  \item El \textbf{Capítulo~\ref{cap:objetivos}} formula los objetivos y la
        metodología general del trabajo.
  \item El \textbf{Capítulo~\ref{cap:estado-arte}} revisa el estado del arte en
        análisis de mamografías y en TDA aplicado a imagen médica.
  \item El \textbf{Capítulo~\ref{cap:fundamentos}} introduce los fundamentos
        teóricos de la homología persistente y del aprendizaje profundo.
  \item El \textbf{Capítulo~\ref{cap:metodologia}} describe la metodología
        propuesta: datos, preprocesado, extracción de descriptores y modelos.
  \item El \textbf{Capítulo~\ref{cap:implementacion}} detalla la implementación
        y las herramientas utilizadas.
  \item El \textbf{Capítulo~\ref{cap:resultados}} presenta y discute los
        resultados experimentales.
  \item El \textbf{Capítulo~\ref{cap:conclusiones}} recoge las conclusiones y
        las líneas de trabajo futuro.
\end{itemize}

\chapter{Objetivos y metodología}
\label{cap:objetivos}

\section{Objetivo general}
El objetivo general de este Trabajo de Fin de Grado es \textbf{investigar la
aplicación del Análisis Topológico de Datos (TDA) a la clasificación y
visualización de mamografías digitales}, evaluando su capacidad discriminativa
entre lesiones benignas y malignas y su complementariedad con los modelos de
aprendizaje profundo.

\section{Objetivos específicos}
Para alcanzar el objetivo general se plantean los siguientes objetivos
específicos:

\begin{enumerate}
  \item[\textbf{O1.}] \textbf{Revisar el estado del arte} en (i) análisis
        automático de mamografías y (ii) aplicaciones del TDA en imagen médica.
  \item[\textbf{O2.}] \textbf{Preparar y preprocesar} el conjunto de datos
        \mbox{CBIS-DDSM}: descarga, limpieza de metadatos, extracción de ROIs y
        normalización de las imágenes.
  \item[\textbf{O3.}] \textbf{Extraer descriptores topológicos} (diagramas de
        persistencia, códigos de barras e imágenes de persistencia) a partir de
        las mamografías mediante homología persistente.
  \item[\textbf{O4.}] \textbf{Entrenar y evaluar clasificadores} sobre los
        descriptores topológicos (p.\,ej. SVM, \textit{random forest}) para la
        tarea benigno/maligno.
  \item[\textbf{O5.}] \textbf{Comparar} el enfoque topológico con un modelo de
        \textit{deep learning} de referencia (CNN con \textit{transfer
        learning}) y \textbf{explorar su combinación} (fusión de
        características).
  \item[\textbf{O6.}] \textbf{Desarrollar herramientas de visualización} que
        hagan interpretables los resultados topológicos para su uso como apoyo
        al diagnóstico.
\end{enumerate}

\section{Metodología general}
El desarrollo del trabajo sigue una metodología incremental e iterativa,
estructurada en las siguientes fases, alineadas con los objetivos anteriores:

\begin{enumerate}
  \item \textbf{Fase de documentación} (O1): estudio de la bibliografía y de las
        herramientas disponibles.
  \item \textbf{Fase de datos} (O2): construcción de un \textit{pipeline}
        reproducible de preparación de datos.
  \item \textbf{Fase de extracción de características} (O3): implementación de la
        extracción de descriptores topológicos.
  \item \textbf{Fase de modelado} (O4, O5): entrenamiento, validación y
        comparación de modelos, evitando la fuga de datos (\textit{data
        leakage}) mediante una separación estricta \textit{train/validation/test}.
  \item \textbf{Fase de análisis y visualización} (O6): interpretación de los
        resultados y desarrollo de visualizaciones.
  \item \textbf{Fase de redacción} de la memoria.
\end{enumerate}

\begin{table}[H]
\centering
\caption{Trazabilidad entre objetivos específicos y capítulos de la memoria.}
\label{tab:trazabilidad}
\begin{tabular}{@{}lll@{}}
\toprule
\textbf{Objetivo} & \textbf{Descripción breve} & \textbf{Capítulo} \\
\midrule
O1 & Estado del arte            & \ref{cap:estado-arte} \\
O2 & Preparación de datos       & \ref{cap:metodologia}, \ref{cap:implementacion} \\
O3 & Descriptores topológicos   & \ref{cap:fundamentos}, \ref{cap:metodologia} \\
O4 & Clasificadores sobre TDA   & \ref{cap:metodologia}, \ref{cap:resultados} \\
O5 & Comparación con DL         & \ref{cap:resultados} \\
O6 & Visualización              & \ref{cap:implementacion}, \ref{cap:resultados} \\
\bottomrule
\end{tabular}
\end{table}

\chapter{Estado del arte}
\label{cap:estado-arte}

Este capítulo revisa los trabajos previos relevantes en dos líneas
complementarias: el análisis automático de mamografías mediante aprendizaje
profundo, y las aplicaciones del Análisis Topológico de Datos en imagen médica.

\section{Análisis automático de mamografías}
El análisis asistido por computador de mamografías tiene una larga trayectoria,
desde los primeros sistemas CAD basados en características diseñadas manualmente
(\textit{hand-crafted features}) hasta los actuales enfoques de aprendizaje
profundo. Las redes neuronales convolucionales (CNN) se han consolidado como el
paradigma dominante en visión por computador y en imagen médica, gracias a su
capacidad para aprender jerarquías de características directamente de los
datos~\cite{litjens2017survey, alieli2024survey}.

En el ámbito específico de la mamografía, destacan trabajos orientados a la
detección y clasificación de masas y microcalcificaciones, así como a la
estimación del riesgo. Yala \textit{et al.}~\cite{yala2019} propusieron un
modelo basado en \textit{deep learning} que mejora la predicción del riesgo de
cáncer de mama a partir de la mamografía, superando a los modelos de riesgo
tradicionales. El uso de \textit{transfer learning} a partir de redes
preentrenadas en ImageNet (ResNet, DenseNet, EfficientNet, etc.) es una práctica
habitual para mitigar la escasez de datos etiquetados en el dominio médico.

\subsection{Conjuntos de datos públicos}
Entre los conjuntos de datos públicos de mamografía más utilizados destaca
\mbox{CBIS-DDSM}~\cite{lee2017cbisddsm}, una versión curada y estandarizada del
clásico DDSM que incluye segmentaciones de ROI verificadas y etiquetas de
patología. Otros recursos relevantes son INbreast y las bases de datos de
ecografía mamaria.

\subsection{Limitaciones}
A pesar de sus buenos resultados, los modelos de \textit{deep learning} en
mamografía presentan limitaciones bien conocidas:
\begin{itemize}
  \item \textbf{Interpretabilidad limitada}: su naturaleza de caja negra
        dificulta la confianza clínica y la certificación.
  \item \textbf{Dependencia de grandes volúmenes de datos} etiquetados.
  \item \textbf{Riesgo de fuga de datos} (\textit{data leakage}) y de sesgos
        de adquisición, que pueden producir métricas artificialmente altas si el
        protocolo experimental no es riguroso.
\end{itemize}
Estas limitaciones motivan la búsqueda de representaciones más robustas e
interpretables, como las que proporciona el TDA.

\section{Análisis Topológico de Datos en imagen médica}
El Análisis Topológico de Datos (TDA) engloba un conjunto de técnicas que
estudian la forma de los datos mediante herramientas de topología
algebraica~\cite{carlsson2009topology, edelsbrunner2010computational}. Su
herramienta central, la \emph{homología persistente}, ha sido aplicada con éxito
en diversos problemas de imagen médica: caracterización de texturas de tumores,
análisis de imágenes histopatológicas, neuroimagen y detección de estructuras en
radiología.

En el contexto de la mamografía y del cáncer de mama, los descriptores
topológicos permiten capturar patrones de conectividad y estructura (número de
componentes conexas, agujeros, ramificaciones de conductos, distribución de
microcalcificaciones) que resultan difíciles de modelar con características
convencionales. Su robustez frente al ruido y su interpretabilidad geométrica
los convierten en un complemento natural de las CNN.

\section{Síntesis y aportación de este trabajo}
La revisión anterior muestra que, si bien el \textit{deep learning} domina el
análisis de mamografías, el TDA ofrece una vía prometedora y poco explorada para
obtener representaciones interpretables y robustas. La aportación de este TFG es
un estudio sistemático de descriptores topológicos sobre \mbox{CBIS-DDSM}, su
evaluación como características para la clasificación benigno/maligno, y su
comparación y combinación con un modelo de \textit{deep learning} de referencia.

\chapter{Fundamentos teóricos}
\label{cap:fundamentos}

Este capítulo introduce los conceptos necesarios para comprender el trabajo: los
fundamentos del Análisis Topológico de Datos, con énfasis en la homología
persistente, y una breve descripción de los modelos de aprendizaje profundo
empleados como referencia.

\section{Análisis Topológico de Datos}
El TDA parte de la idea de que los datos poseen una \emph{forma} y que dicha
forma contiene información relevante. Formalmente, se construye a partir de los
datos una familia de espacios topológicos y se estudia cómo evolucionan sus
invariantes.

\subsection{Complejos simpliciales y filtraciones}
\begin{definicion}[Complejo simplicial]
Un \emph{complejo simplicial} $K$ es un conjunto de símplices (vértices, aristas,
triángulos, tetraedros, \dots) cerrado bajo la operación de tomar caras y tal
que la intersección de dos símplices es una cara de ambos.
\end{definicion}

Para analizar datos se construye una \emph{filtración}, es decir, una familia
anidada de complejos simpliciales
\[
  \emptyset = K_0 \subseteq K_1 \subseteq \cdots \subseteq K_n = K,
\]
indexada por un parámetro de escala. En imágenes, una filtración habitual es la
\emph{filtración por sublevel sets} de la intensidad: se van incluyendo los
píxeles cuyo valor es menor o igual que un umbral creciente.

\subsection{Homología y números de Betti}
La homología asocia a cada complejo una secuencia de grupos cuyos rangos, los
\emph{números de Betti} $\beta_k$, cuentan las características topológicas de
dimensión $k$:
\begin{itemize}
  \item $\beta_0$: número de \textbf{componentes conexas}.
  \item $\beta_1$: número de \textbf{agujeros} o ciclos unidimensionales.
  \item $\beta_2$: número de \textbf{cavidades} o huecos tridimensionales.
\end{itemize}

\subsection{Homología persistente}
La \emph{homología persistente} rastrea el nacimiento (\textit{birth}) y la
muerte (\textit{death}) de las características topológicas a lo largo de la
filtración. Una característica que persiste durante un amplio rango de escalas se
considera una propiedad estructural significativa, mientras que las de vida corta
suelen corresponder a ruido.

\begin{definicion}[Diagrama de persistencia]
El \emph{diagrama de persistencia} es un multiconjunto de puntos
$(b_i, d_i) \in \mathbb{R}^2$ con $b_i \le d_i$, donde cada punto representa una
característica topológica que nace en la escala $b_i$ y muere en $d_i$. La
\emph{persistencia} de la característica es $d_i - b_i$.
\end{definicion}

La misma información puede representarse como un \emph{código de barras}
(\textit{barcode}), donde cada característica es un intervalo $[b_i, d_i]$.

\subsection{Estabilidad y vectorización}
Una propiedad fundamental es la \textbf{estabilidad}: pequeñas perturbaciones en
los datos producen pequeñas variaciones en el diagrama de persistencia (medidas
con la distancia \textit{bottleneck} o de Wasserstein). Esto explica la robustez
del TDA frente al ruido.

Para poder utilizar los diagramas como entrada de algoritmos de aprendizaje
automático es necesario \emph{vectorizarlos}. Entre las representaciones más
utilizadas destacan:
\begin{itemize}
  \item \textbf{Imágenes de persistencia} (\textit{persistence images}):
        transforman el diagrama en una imagen (matriz) de tamaño fijo.
  \item \textbf{Persistence landscapes}: funciones a trozos que resumen el
        diagrama.
  \item \textbf{Curvas de Betti} y estadísticos de persistencia.
\end{itemize}

\section{Aprendizaje profundo (referencia)}
Como referencia comparativa se emplean redes neuronales convolucionales (CNN).

\subsection{Redes neuronales convolucionales}
Una CNN aprende, mediante capas convolucionales, jerarquías de características
espaciales a partir de las imágenes. Combinada con capas de \textit{pooling} y
capas totalmente conectadas, permite abordar tareas de clasificación de imagen.

\subsection{Transfer learning}
Dado el tamaño limitado de los conjuntos de datos médicos, se recurre al
\textit{transfer learning}: se parte de una red preentrenada en un gran corpus
(p.\,ej. ImageNet) y se ajusta (\textit{fine-tuning}) a la tarea de mamografía.
En este trabajo se utiliza como referencia una arquitectura de la familia
DenseNet/ResNet.

\begin{ejemplo}[Intuición del enfoque topológico en mamografía]
Una agrupación de microcalcificaciones puede caracterizarse por el número de
componentes conexas ($\beta_0$) y su patrón de aparición a lo largo de la
filtración de intensidad, mientras que la estructura de una masa espiculada
puede reflejarse en la aparición y desaparición de ciclos ($\beta_1$). El TDA
codifica estos patrones de forma compacta y robusta.
\end{ejemplo}

\chapter{Metodología propuesta}
\label{cap:metodologia}

En este capítulo se describe la metodología propuesta para aplicar el TDA a la
clasificación de mamografías, desde la preparación de los datos hasta la
evaluación de los modelos.

\section{Visión general del \textit{pipeline}}
El sistema propuesto consta de cinco etapas:
\begin{enumerate}
  \item Preparación de datos (CBIS-DDSM).
  \item Preprocesado de imágenes y extracción de ROIs.
  \item Extracción de descriptores topológicos.
  \item Clasificación (topológica, \textit{deep learning} y fusión).
  \item Evaluación y visualización.
\end{enumerate}

\begin{figure}[H]
\centering
\resizebox{\textwidth}{!}{%
\begin{tikzpicture}[
    node distance = 0.9cm and 1.1cm,
    stage/.style = {
      draw, rounded corners, align=center, font=\small,
      minimum height=1.1cm, text width=2.2cm, fill=blue!6
    },
    arr/.style = {-Latex, thick}
  ]
  \node[stage] (datos)  {Datos\\[-2pt]\footnotesize CBIS-DDSM};
  \node[stage, right=of datos]  (prep)  {Preprocesado\\[-2pt]\footnotesize ROI, CLAHE};
  \node[stage, right=of prep]   (topo)  {Homología\\persistente};
  \node[stage, right=of topo]   (vec)   {Vectorización\\[-2pt]\footnotesize imagen de persist.};
  \node[stage, right=of vec]    (clf)   {Clasificador\\[-2pt]\footnotesize SVM/RF/CNN};
  \node[stage, right=of clf]    (eval)  {Evaluación\\[-2pt]\footnotesize métricas clínicas};

  \draw[arr] (datos) -- (prep);
  \draw[arr] (prep)  -- (topo);
  \draw[arr] (topo)  -- (vec);
  \draw[arr] (vec)   -- (clf);
  \draw[arr] (clf)   -- (eval);
\end{tikzpicture}%
}
\caption{Esquema general del \textit{pipeline} propuesto: de los datos crudos a
las métricas clínicas, pasando por la extracción de descriptores topológicos.}
\label{fig:pipeline}
\end{figure}

\section{Conjunto de datos}
Se utiliza \mbox{CBIS-DDSM}~\cite{lee2017cbisddsm}. El fichero de metadatos
disponible en el repositorio (\texttt{CBIS-DDSM-All-...\,.xlsx}) contiene la
información de cada estudio: identificador de paciente, tipo de lesión (masa o
calcificación), densidad de la mama (BI-RADS), vista (CC/MLO), y la etiqueta de
patología (\texttt{BENIGN}, \texttt{MALIGNANT}, \texttt{BENIGN\_WITHOUT\_
CALLBACK}).

\begin{itemize}
  \item \textbf{Tarea principal}: clasificación binaria benigno vs.\ maligno.
  \item \textbf{Unidad de análisis}: región de interés (ROI) de la lesión y/o
        imagen completa reescalada.
\end{itemize}

\subsection{Protocolo de evaluación y prevención de fuga de datos}
Se adopta una partición estricta en \textbf{entrenamiento / validación / test}.
Para evitar la fuga de datos (\textit{data leakage}), la separación se realiza
\textbf{a nivel de paciente}: todas las imágenes de un mismo paciente pertenecen
a una única partición. Se emplea estratificación por la etiqueta de patología
para mantener la proporción de clases.

\begin{quote}
\textit{Nota metodológica.} En el material de partida
(\texttt{Demo\_Notebook.ipynb}) se observaba una exactitud de entrenamiento del
100\,\% desde la primera época y una exactitud de validación cercana al 100\,\%,
lo que es un indicio claro de fuga de datos o de un problema en el
\textit{pipeline}. Este trabajo corrige dicho protocolo mediante partición por
paciente, control de la normalización y validación cruzada.
\end{quote}

\section{Preprocesado de imágenes}
El preprocesado incluye:
\begin{itemize}
  \item Conversión y normalización de intensidades.
  \item Recorte de la ROI a partir de las máscaras de segmentación.
  \item Reescalado a un tamaño común (p.\,ej. $224\times224$).
  \item Opcionalmente, realce de contraste (CLAHE) y filtrado de ruido.
\end{itemize}

\section{Extracción de descriptores topológicos}
A partir de cada imagen (o ROI) se calcula la homología persistente mediante una
filtración por \textit{sublevel sets} de la intensidad (para imágenes en escala
de grises) o sobre nubes de puntos derivadas de las microcalcificaciones. Se
obtienen los diagramas de persistencia en dimensiones $0$ y $1$, y a partir de
ellos se generan los vectores de características:
\begin{itemize}
  \item Imágenes de persistencia.
  \item Estadísticos de persistencia (número de puntos, persistencia total,
        máxima, media; curvas de Betti).
\end{itemize}

\section{Modelos de clasificación}
Se comparan tres enfoques:
\begin{enumerate}
  \item \textbf{Topológico}: clasificadores clásicos (SVM, \textit{random
        forest}, regresión logística) sobre los descriptores topológicos.
  \item \textbf{Deep learning}: CNN con \textit{transfer learning} sobre las
        imágenes.
  \item \textbf{Fusión}: combinación de las características topológicas con las
        \textit{embeddings} de la CNN.
\end{enumerate}

\section{Métricas de evaluación}
Dado el carácter clínico del problema, se emplean varias métricas
complementarias:
\begin{itemize}
  \item \textbf{Exactitud} (\textit{accuracy}).
  \item \textbf{Sensibilidad} (\textit{recall}) --- crítica para no pasar por
        alto casos malignos.
  \item \textbf{Especificidad} y \textbf{precisión}.
  \item \textbf{F1-score} y \textbf{AUC-ROC}.
  \item \textbf{Matriz de confusión}.
\end{itemize}
La sensibilidad recibe especial atención, ya que un falso negativo (no detectar
un cáncer) tiene un coste clínico mucho mayor que un falso positivo.

\chapter{Implementación}
\label{cap:implementacion}

Este capítulo describe las herramientas, la organización del código y los
detalles de implementación del \textit{pipeline}.

\section{Herramientas y librerías}
El proyecto se desarrolla en \textbf{Python}. Las principales librerías son:

\begin{table}[H]
\centering
\caption{Principales herramientas utilizadas.}
\label{tab:herramientas}
\begin{tabular}{@{}ll@{}}
\toprule
\textbf{Ámbito} & \textbf{Herramienta} \\
\midrule
Manipulación de datos      & pandas, NumPy \\
Imagen                     & OpenCV, scikit-image, Pillow \\
TDA / homología persistente& giotto-tda, Ripser, GUDHI \\
Aprendizaje automático     & scikit-learn \\
Aprendizaje profundo       & TensorFlow / Keras (o PyTorch) \\
Visualización              & Matplotlib, seaborn \\
Entorno                    & Jupyter Notebook \\
\bottomrule
\end{tabular}
\end{table}

\section{Organización del código}
El repositorio se estructura de la siguiente forma (véase el
\texttt{README.md}):
\begin{lstlisting}[language=bash, numbers=none]
mammoviz-tda/
  data/            # datos (no versionados)
  src/             # codigo fuente del pipeline
    data.py        # carga y particion por paciente
    preprocessing.py  # preprocesado de imagenes
    topology.py    # homologia persistente y vectorizacion
    models.py      # clasificadores y CNN
    evaluate.py    # metricas y visualizacion
  notebooks/       # experimentos reproducibles
  tfg/             # memoria en LaTeX
\end{lstlisting}

\section{Extracción de descriptores topológicos}
El siguiente fragmento ilustra el cálculo de un diagrama de persistencia a partir
de una imagen en escala de grises mediante una filtración por \textit{sublevel
sets}, utilizando \texttt{giotto-tda}.

\begin{lstlisting}[language=Python, caption={Ejemplo de extracción de un diagrama de persistencia.}, label={lst:tda}]
import numpy as np
from gtda.images import HeightFiltration, RadialFiltration
from gtda.homology import CubicalPersistence
from gtda.diagrams import PersistenceImage

def descriptor_topologico(imagen_gris):
    # imagen_gris: array 2D normalizado en [0, 1]
    x = imagen_gris[None, :, :]              # (1, H, W)
    cubical = CubicalPersistence(homology_dimensions=(0, 1))
    diagramas = cubical.fit_transform(x)     # diagramas de persistencia
    pi = PersistenceImage(n_bins=20)
    vector = pi.fit_transform(diagramas)     # imagen de persistencia
    return vector.reshape(len(x), -1)
\end{lstlisting}

\section{Reproducibilidad}
Para garantizar la reproducibilidad de los experimentos:
\begin{itemize}
  \item Se fija una semilla aleatoria común.
  \item Se documentan las versiones de las librerías (\texttt{requirements.txt}).
  \item La partición de datos por paciente se guarda en disco para reutilizarla.
  \item Cada experimento se ejecuta desde un \textit{notebook} autocontenido.
\end{itemize}

\section{Visualización}
Se desarrollan visualizaciones específicas para el análisis topológico:
diagramas de persistencia y códigos de barras superpuestos sobre la imagen,
mapas de calor de las imágenes de persistencia, y comparación de las regiones
más informativas frente a las activaciones de la CNN (p.\,ej. Grad-CAM).

\chapter{Experimentos y resultados}
\label{cap:resultados}

Este capítulo presenta el diseño experimental y los resultados obtenidos.
\textit{(Este capítulo debe completarse con los resultados reales a medida que
se ejecuten los experimentos; a continuación se proporciona la estructura y las
plantillas de tablas y figuras.)}

\section{Diseño experimental}
Se definen los siguientes experimentos, alineados con los objetivos O4 y O5:
\begin{description}
  \item[E1.] Clasificación con descriptores topológicos + clasificador clásico.
  \item[E2.] Clasificación con CNN (\textit{transfer learning}).
  \item[E3.] Fusión de características topológicas y \textit{embeddings} de CNN.
  \item[E4.] Análisis de robustez frente a ruido y a variaciones de adquisición.
\end{description}

\section{Configuración}
Se detallan los hiperparámetros, la partición de datos por paciente, el número
de repeticiones y el hardware utilizado.

\section{Resultados}

\begin{table}[H]
\centering
\caption{Resultados comparativos en el conjunto de test (plantilla).}
\label{tab:resultados}
\begin{tabular}{@{}lccccc@{}}
\toprule
\textbf{Modelo} & \textbf{Exactitud} & \textbf{Sensib.} & \textbf{Especif.} &
\textbf{F1} & \textbf{AUC} \\
\midrule
Baseline (clase mayoritaria) & -- & -- & -- & -- & -- \\
TDA + SVM                    & -- & -- & -- & -- & -- \\
TDA + Random Forest          & -- & -- & -- & -- & -- \\
CNN (transfer learning)      & -- & -- & -- & -- & -- \\
Fusión TDA + CNN             & -- & -- & -- & -- & -- \\
\bottomrule
\end{tabular}
\end{table}

\begin{figure}[H]
\centering
% \includegraphics[width=0.7\textwidth]{figuras/roc.png}
\fbox{\parbox[c][4cm][c]{0.7\textwidth}{\centering\itshape
Curvas ROC de los distintos modelos. (Sustituir por la figura definitiva.)}}
\caption{Comparación de curvas ROC.}
\label{fig:roc}
\end{figure}

\section{Discusión}
Se analizan los resultados a la luz de los objetivos: capacidad discriminativa de
los descriptores topológicos, comparación con el \textit{deep learning},
beneficio de la fusión, robustez frente al ruido e interpretabilidad. Se discuten
también las limitaciones observadas.

\chapter{Conclusiones y trabajo futuro}
\label{cap:conclusiones}

\section{Conclusiones}
\textit{(Redactar a partir de los resultados obtenidos.)} En este Trabajo de Fin
de Grado se ha estudiado la aplicación del Análisis Topológico de Datos a la
clasificación y visualización de mamografías digitales del conjunto
\mbox{CBIS-DDSM}. Las principales conclusiones, en relación con los objetivos
planteados, son:

\begin{itemize}
  \item Se ha revisado el estado del arte en análisis de mamografías y en TDA
        aplicado a imagen médica (O1).
  \item Se ha construido un \textit{pipeline} reproducible de preparación de
        datos con partición por paciente que evita la fuga de datos (O2).
  \item Se han extraído y evaluado descriptores topológicos como características
        para la clasificación benigno/maligno (O3, O4).
  \item Se ha comparado el enfoque topológico con un modelo de \textit{deep
        learning} de referencia y se ha explorado su combinación (O5).
  \item Se han desarrollado visualizaciones que mejoran la interpretabilidad de
        los resultados (O6).
\end{itemize}

\section{Grado de cumplimiento de los objetivos}
Se valora el grado de cumplimiento de cada objetivo específico y las
desviaciones respecto a la planificación inicial.

\section{Limitaciones}
Se discuten las limitaciones del trabajo: tamaño y sesgos del conjunto de datos,
coste computacional de la homología persistente, y alcance de la validación
clínica.

\section{Trabajo futuro}
Como líneas de trabajo futuro se proponen:
\begin{itemize}
  \item Extender el estudio a otras filtraciones y descriptores topológicos
        (\textit{multiparameter persistence}).
  \item Integrar el TDA directamente en arquitecturas de \textit{deep learning}
        (capas topológicas diferenciables).
  \item Validar el enfoque en otros conjuntos de datos (INbreast, datos propios).
  \item Realizar una evaluación con radiólogos para medir la utilidad clínica de
        las visualizaciones.
\end{itemize}


%------------------------------- Bibliografía ---------------------------------
\printbibliography[heading=bibintoc,title={Bibliografía}]

%------------------------------- Apéndices ------------------------------------
\appendix
\chapter{Material adicional}
\label{ap:material}

\section{Instalación del entorno}
Instrucciones para reproducir el entorno de desarrollo:
\begin{lstlisting}[language=bash, numbers=none]
python -m venv .venv
source .venv/bin/activate        # En Windows: .venv\Scripts\activate
pip install -r requirements.txt
\end{lstlisting}

\section{Descripción de los campos de CBIS-DDSM}
Tabla de referencia con los campos del fichero de metadatos y su significado
(densidad BI-RADS, tipo de lesión, vista, patología, etc.).

\section{Fragmentos de código adicionales}
Fragmentos de código relevantes que, por extensión, no se incluyen en el cuerpo
de la memoria.


\end{document}
